\documentclass[12pt,a4paper]{article}

% Vietnamese support
\usepackage[utf8]{inputenc}
\usepackage[vietnamese]{babel}
\usepackage{vntex}

% Essential packages
\usepackage{graphicx}
\usepackage{amsmath}
\usepackage{amssymb}
\usepackage{booktabs}
\usepackage{longtable}
\usepackage{array}
\usepackage{multirow}
\usepackage{hyperref}
\usepackage{xcolor}
\usepackage{listings}
\usepackage{fancyvrb}
\usepackage{geometry}
\usepackage{titlesec}
\usepackage{tocloft}
\usepackage{fancyhdr}
\usepackage{tikz}
\usetikzlibrary{shapes,arrows,positioning,calc}

% Page geometry
\geometry{
    left=2.5cm,
    right=2cm,
    top=2.5cm,
    bottom=2.5cm
}

% Colors
\definecolor{codegreen}{rgb}{0,0.6,0}
\definecolor{codegray}{rgb}{0.5,0.5,0.5}
\definecolor{codepurple}{rgb}{0.58,0,0.82}
\definecolor{backcolour}{rgb}{0.95,0.95,0.92}
\definecolor{darkblue}{rgb}{0.0,0.0,0.6}

% Code listing style
\lstdefinestyle{mystyle}{
    backgroundcolor=\color{backcolour},
    commentstyle=\color{codegreen},
    keywordstyle=\color{magenta},
    numberstyle=\tiny\color{codegray},
    stringstyle=\color{codepurple},
    basicstyle=\ttfamily\footnotesize,
    breakatwhitespace=false,
    breaklines=true,
    captionpos=b,
    keepspaces=true,
    numbers=left,
    numbersep=5pt,
    showspaces=false,
    showstringspaces=false,
    showtabs=false,
    tabsize=4,
    frame=single,
    rulecolor=\color{black}
}
\lstset{style=mystyle}

% C# language definition
\lstdefinelanguage{CSharp}{
    keywords={abstract,as,base,bool,break,byte,case,catch,char,checked,class,const,continue,decimal,default,delegate,do,double,else,enum,event,explicit,extern,false,finally,fixed,float,for,foreach,goto,if,implicit,in,int,interface,internal,is,lock,long,namespace,new,null,object,operator,out,override,params,private,protected,public,readonly,ref,return,sbyte,sealed,short,sizeof,stackalloc,static,string,struct,switch,this,throw,true,try,typeof,uint,ulong,unchecked,unsafe,ushort,using,virtual,void,volatile,while,async,await,var,dynamic},
    sensitive=true,
    morecomment=[l]{//},
    morecomment=[s]{/*}{*/},
    morestring=[b]",
    morestring=[b]'
}

% Hyperref setup
\hypersetup{
    colorlinks=true,
    linkcolor=darkblue,
    filecolor=magenta,
    urlcolor=blue,
    citecolor=darkblue
}

% Header and footer
\pagestyle{fancy}
\fancyhf{}
\rhead{Session 5 - Parallel Programming}
\lhead{Reduction Pattern}
\cfoot{\thepage}

% Title formatting
\titleformat{\section}
    {\normalfont\Large\bfseries}{\thesection.}{1em}{}
\titleformat{\subsection}
    {\normalfont\large\bfseries}{\thesubsection.}{1em}{}
\titleformat{\subsubsection}
    {\normalfont\normalsize\bfseries}{\thesubsubsection.}{1em}{}

% Document info
\title{
    \vspace{-1cm}
    \textbf{Báo cáo: Lập trình Song song trong C\#}\\
    \Large Reduction Pattern\\
    \vspace{0.5cm}
    \large Session 5: Parallel Programming \& Reduction Techniques
}
\author{}
\date{Ngày 03/02/2026}

\begin{document}

\maketitle
\thispagestyle{empty}

\vspace{1cm}

\tableofcontents
\newpage

%==============================================================================
\section{Giới thiệu}
%==============================================================================

\subsection{Tổng quan về Lập trình Song song}

Lập trình song song (Parallel Programming) là kỹ thuật chia nhỏ một tác vụ lớn thành nhiều tác vụ nhỏ hơn để thực thi đồng thời trên nhiều nhân CPU, từ đó tăng hiệu suất và giảm thời gian xử lý.

\subsection{Reduction Pattern}

\textbf{Reduction} là một design pattern phổ biến trong lập trình song song, dùng để kết hợp (combine) nhiều giá trị thành một giá trị duy nhất thông qua một phép toán kết hợp (associative operation) như:

\begin{itemize}
    \item Tính tổng (Sum)
    \item Tính tích (Product)
    \item Tìm min/max
    \item Ghép nối (Concatenation)
\end{itemize}

\subsection{Các công cụ được sử dụng trong C\#}

\begin{table}[h]
\centering
\begin{tabular}{|l|l|l|}
\hline
\textbf{Công cụ} & \textbf{Mô tả} & \textbf{Use case} \\
\hline
\texttt{Parallel.For} & Vòng lặp song song & Tính toán độc lập trên từng phần tử \\
\hline
\texttt{Channel<T>} & Hàng đợi thread-safe & Producer-Consumer pattern \\
\hline
\texttt{Interlocked} & Atomic operations & Cập nhật biến chia sẻ an toàn \\
\hline
\texttt{PLINQ} & Parallel LINQ & Truy vấn song song trên collections \\
\hline
\end{tabular}
\caption{Các công cụ lập trình song song trong C\#}
\end{table}

%==============================================================================
\section{Bài toán Factorial}
%==============================================================================

\subsection{Mô tả bài toán}

Tính giai thừa của các số từ 0 đến n-1. Bài toán này minh họa hai cách tiếp cận song song cơ bản.

\subsection{Cách tiếp cận 1: Sử dụng Parallel.For}

\subsubsection{Ý tưởng}

\begin{itemize}
    \item Mỗi phép tính $i!$ là \textbf{độc lập} với các phép tính khác
    \item Sử dụng \texttt{Parallel.For} để phân phối công việc tự động cho các thread
\end{itemize}

\subsubsection{Sơ đồ thực thi}

\begin{figure}[h]
\centering
\begin{tikzpicture}[
    box/.style={rectangle, draw, minimum width=2cm, minimum height=1.5cm, align=center},
    arrow/.style={->, >=stealth, thick}
]
    % Main box
    \node[rectangle, draw, minimum width=12cm, minimum height=0.8cm, fill=gray!20] (main) at (0,2) {\texttt{Parallel.For(0, n)}};
    
    % Thread boxes
    \node[box, fill=blue!10] (t1) at (-4.5,0) {Thread 1\\$0! = 1$\\$7! = 5040$};
    \node[box, fill=green!10] (t2) at (-1.5,0) {Thread 2\\$1! = 1$\\$4! = 24$};
    \node[box, fill=yellow!10] (t3) at (1.5,0) {Thread 3\\$2! = 2$\\$5! = 120$};
    \node[box, fill=red!10] (t4) at (4.5,0) {Thread 4\\$3! = 6$\\$6! = 720$};
    
    % Note
    \node[align=center] at (0,-1.5) {\textit{Thứ tự output KHÔNG được đảm bảo}};
\end{tikzpicture}
\caption{Sơ đồ thực thi Parallel.For}
\end{figure}

\subsubsection{Code minh họa}

\begin{lstlisting}[language=CSharp, caption={Factorial với Parallel.For}]
Parallel.For(0, n, i =>
{
    long factorial = ComputeFactorial(i);
    Console.WriteLine($"{i}! = {factorial}");
});
\end{lstlisting}

\subsubsection{Kết quả thực thi}

\begin{Verbatim}[frame=single, fontsize=\small]
=== Factorial using Parallel.For ===

Computing factorials from 0! to 9!:

0! = 1
7! = 5040
2! = 2
1! = 1
9! = 362880
4! = 24
6! = 720
5! = 120
8! = 40320
3! = 6

--- Parallel execution completed ---
Note: Results may appear in different order due to parallel execution
\end{Verbatim}

\textbf{Quan sát:} Thứ tự in ra không theo thứ tự tăng dần (0, 7, 2, 1, 9...) vì các thread chạy song song và hoàn thành tại các thời điểm khác nhau.

\subsection{Cách tiếp cận 2: Sử dụng Channel (Producer-Consumer)}

\subsubsection{Ý tưởng}

\begin{itemize}
    \item \textbf{Producer}: Tính toán giai thừa và gửi kết quả vào channel
    \item \textbf{Consumer}: Đọc kết quả từ channel và hiển thị
    \item Channel đảm bảo \textbf{thread-safe} và \textbf{thứ tự} (FIFO)
\end{itemize}

\subsubsection{Sơ đồ Producer-Consumer}

\begin{figure}[h]
\centering
\begin{tikzpicture}[
    box/.style={rectangle, draw, minimum width=3cm, minimum height=3cm, align=center},
    arrow/.style={->, >=stealth, thick}
]
    % Producer
    \node[box, fill=blue!15] (producer) at (-4,0) {\textbf{PRODUCER}\\[0.3cm]Compute 0!\\Compute 1!\\...\\Compute 9!\\Complete()};
    
    % Channel
    \node[box, fill=yellow!15] (channel) at (0,0) {\textbf{CHANNEL}\\[0.3cm](Buffer)\\FIFO};
    
    % Consumer
    \node[box, fill=green!15] (consumer) at (4,0) {\textbf{CONSUMER}\\[0.3cm]Display 0!\\Display 1!\\...\\Display 9!\\End};
    
    % Arrows
    \draw[arrow] (producer) -- (channel);
    \draw[arrow] (channel) -- (consumer);
\end{tikzpicture}
\caption{Sơ đồ Producer-Consumer Pattern}
\end{figure}

\subsubsection{Code minh họa}

\begin{lstlisting}[language=CSharp, caption={Factorial với Channel}]
// Tao channel
var channel = Channel.CreateUnbounded<(int number, long factorial)>();
var writer = channel.Writer;
var reader = channel.Reader;

// Producer
var producerTask = Task.Run(async () =>
{
    for (int i = 0; i < n; i++)
    {
        long factorial = ComputeFactorial(i);
        await writer.WriteAsync((i, factorial));
    }
    writer.Complete();
});

// Consumer
await foreach (var item in reader.ReadAllAsync())
{
    Console.WriteLine($"[Consumer] Received: {item.number}! = {item.factorial}");
}
\end{lstlisting}

\subsubsection{Kết quả thực thi}

\begin{Verbatim}[frame=single, fontsize=\small]
=== Factorial using Channel (Producer-Consumer) ===

[Producer] Starting to compute factorials...

[Producer] Computed and sent: 0! = 1
[Producer] Computed and sent: 1! = 1
...
[Producer] Computed and sent: 9! = 362880

[Producer] Completed writing all factorials.

[Consumer] Received: 0! = 1
[Consumer] Received: 1! = 1
...
[Consumer] Received: 9! = 362880

[Consumer] Finished reading all results.
\end{Verbatim}

\textbf{Quan sát:} Kết quả được Consumer nhận theo đúng thứ tự FIFO do Channel đảm bảo.

\subsection{So sánh hai cách tiếp cận}

\begin{table}[h]
\centering
\begin{tabular}{|l|c|c|}
\hline
\textbf{Tiêu chí} & \textbf{Parallel.For} & \textbf{Channel} \\
\hline
Độ phức tạp & Đơn giản & Phức tạp hơn \\
\hline
Thứ tự kết quả & Không đảm bảo & Đảm bảo FIFO \\
\hline
Mô hình & Data parallelism & Producer-Consumer \\
\hline
Use case & Tính toán độc lập & Pipeline processing \\
\hline
\end{tabular}
\caption{So sánh Parallel.For và Channel}
\end{table}

%==============================================================================
\section{Bài toán Vector Sum}
%==============================================================================

\subsection{Mô tả bài toán}

Tính tổng tất cả phần tử trong một vector (mảng) có kích thước lớn. Đây là bài toán \textbf{Reduction} điển hình.

\begin{equation}
\text{Vector} = [1, 2, 3, ..., 1{,}000{,}000]
\end{equation}

\begin{equation}
\text{Sum} = \sum_{i=1}^{1{,}000{,}000} i = \frac{n(n+1)}{2} = 500{,}000{,}500{,}000
\end{equation}

\subsection{Các phương pháp tiếp cận}

\subsubsection{Method 1: Sequential Sum (Tuần tự)}

\begin{lstlisting}[language=CSharp, caption={Sequential Sum}]
long sum = 0;
for (int i = 0; i < vectorSize; i++)
{
    sum += vector[i];
}
\end{lstlisting}

\textbf{Đặc điểm:} Đơn giản, dễ hiểu, nhưng không tận dụng được đa nhân CPU.

\subsubsection{Method 2: Parallel with Interlocked}

\begin{lstlisting}[language=CSharp, caption={Parallel với Interlocked}]
long parallelSum = 0;
Parallel.For(0, vectorSize, i =>
{
    Interlocked.Add(ref parallelSum, vector[i]);
});
\end{lstlisting}

\textbf{Vấn đề:} Sử dụng \texttt{Interlocked.Add} cho mỗi phần tử gây ra \textbf{contention} (tranh chấp) cao giữa các thread.

\subsubsection{Method 3: Parallel with Local State (Khuyến nghị)}

\begin{lstlisting}[language=CSharp, caption={Parallel với Local State}]
long parallelSum = 0;
object lockObj = new object();

Parallel.For(0, vectorSize,
    () => 0L,                              // Init local state
    (i, loopState, localSum) =>            // Body
    {
        return localSum + vector[i];
    },
    (localSum) =>                          // Final (combine)
    {
        lock (lockObj) { parallelSum += localSum; }
    }
);
\end{lstlisting}

\subsubsection{Sơ đồ Reduction với Local State}

\begin{figure}[h]
\centering
\begin{tikzpicture}[
    box/.style={rectangle, draw, minimum width=2.5cm, minimum height=1.8cm, align=center, font=\small},
    phase/.style={rectangle, draw, minimum width=11cm, minimum height=0.7cm, fill=gray!20},
    arrow/.style={->, >=stealth, thick}
]
    % Phase 1 header
    \node[phase] at (0,3) {\textbf{PHASE 1: Parallel Accumulate}};
    
    % Thread boxes
    \node[box, fill=blue!10] (t1) at (-4,1) {Thread 1\\localSum = 0\\+= v[0..249k]\\= partial\_1};
    \node[box, fill=green!10] (t2) at (-1.3,1) {Thread 2\\localSum = 0\\+= v[250k..499k]\\= partial\_2};
    \node[box, fill=yellow!10] (t3) at (1.4,1) {Thread 3\\localSum = 0\\+= v[500k..749k]\\= partial\_3};
    \node[box, fill=red!10] (t4) at (4.1,1) {Thread N\\localSum = 0\\+= v[750k..999k]\\= partial\_N};
    
    % Arrows down
    \draw[arrow] (t1.south) -- ++(0,-0.5);
    \draw[arrow] (t2.south) -- ++(0,-0.5);
    \draw[arrow] (t3.south) -- ++(0,-0.5);
    \draw[arrow] (t4.south) -- ++(0,-0.5);
    
    % Phase 2
    \node[rectangle, draw, minimum width=11cm, minimum height=1.2cm, fill=orange!15] at (0,-1) {\textbf{PHASE 2: Final Combine (lock)}\\totalSum = partial\_1 + partial\_2 + partial\_3 + ... + partial\_N};
\end{tikzpicture}
\caption{Sơ đồ Reduction với Local State}
\end{figure}

\subsubsection{Method 4: Manual Partitioning}

\begin{lstlisting}[language=CSharp, caption={Manual Partitioning}]
int numPartitions = Environment.ProcessorCount;
long[] partialSums = new long[numPartitions];

Parallel.For(0, numPartitions, p =>
{
    int start = p * partitionSize;
    int end = (p == numPartitions - 1) ? vectorSize : start + partitionSize;
    
    long localSum = 0;
    for (int i = start; i < end; i++)
        localSum += vector[i];
    partialSums[p] = localSum;
});

// Combine
for (int p = 0; p < numPartitions; p++)
    parallelSum += partialSums[p];
\end{lstlisting}

\textbf{Đặc điểm:} Kiểm soát hoàn toàn việc phân chia công việc.

\subsubsection{Method 5: PLINQ}

\begin{lstlisting}[language=CSharp, caption={PLINQ}]
long plinqSum = vector.AsParallel().Sum(x => (long)x);
\end{lstlisting}

\textbf{Đặc điểm:} Cú pháp ngắn gọn, declarative, nhưng có overhead.

\subsection{Kết quả thực thi}

\begin{Verbatim}[frame=single, fontsize=\small]
=== Vector Sum using Parallel (Reduction) ===

Vector size: 1,000,000
Vector: [1, 2, 3, ..., 1000000]

--- Method 1: Sequential Sum ---
Sequential Sum = 500,000,500,000
Time: 1 ms

--- Method 2: Parallel with Interlocked ---
Parallel Sum (Interlocked) = 500,000,500,000
Time: 28 ms

--- Method 3: Parallel with Local State ---
Parallel Sum (Local State) = 500,000,500,000
Time: 3 ms

--- Method 4: Parallel with Partitioning ---
Parallel Sum (Partitioning, 12 partitions) = 500,000,500,000
Time: 0 ms

--- Method 5: PLINQ (Parallel LINQ) ---
PLINQ Sum = 500,000,500,000
Time: 8 ms

=== Verification ===
Expected Sum (formula n(n+1)/2) = 500,000,500,000
All results correct: True
\end{Verbatim}

\subsection{Phân tích kết quả}

\begin{table}[h]
\centering
\begin{tabular}{|l|c|l|}
\hline
\textbf{Method} & \textbf{Thời gian} & \textbf{Nhận xét} \\
\hline
Sequential & 1 ms & Baseline, tận dụng CPU cache tốt \\
\hline
Interlocked & 28 ms & \textbf{Chậm nhất} do contention cao \\
\hline
Local State & 3 ms & Hiệu quả, giảm contention \\
\hline
Partitioning & < 1 ms & \textbf{Nhanh nhất}, ít overhead \\
\hline
PLINQ & 8 ms & Tiện lợi nhưng có overhead \\
\hline
\end{tabular}
\caption{So sánh hiệu năng các phương pháp Vector Sum}
\end{table}

\textbf{Quan sát quan trọng:} Với dữ liệu nhỏ (1 triệu phần tử), overhead của parallel có thể làm chậm hơn sequential. Parallel hiệu quả hơn với dữ liệu lớn.

%==============================================================================
\section{Bài toán Array Sum Comparison}
%==============================================================================

\subsection{Mô tả bài toán}

So sánh 3 phương pháp tính tổng mảng với kích thước lớn (100 triệu phần tử):

\begin{enumerate}
    \item \textbf{Sequential} - Tuần tự
    \item \textbf{Parallel} - Song song với Local State
    \item \textbf{Channel} - Producer-Consumer pattern
\end{enumerate}

\subsection{Thiết kế giải pháp}

\begin{figure}[h]
\centering
\begin{tikzpicture}[
    box/.style={rectangle, draw, minimum width=3cm, minimum height=2cm, align=center},
    arrow/.style={->, >=stealth, thick}
]
    % Array
    \node[rectangle, draw, minimum width=12cm, minimum height=0.8cm, fill=gray!20] (array) at (0,3) {ARRAY [1, 2, 3, ..., 100,000,000]};
    
    % Three approaches
    \node[box, fill=blue!15] (seq) at (-4,0.5) {\textbf{SEQUENTIAL}\\Single thread\\Simple loop};
    \node[box, fill=green!15] (par) at (0,0.5) {\textbf{PARALLEL}\\Multi-thread\\Local State};
    \node[box, fill=orange!15] (chan) at (4,0.5) {\textbf{CHANNEL}\\Producer-\\Consumer};
    
    % Arrows
    \draw[arrow] (array.south) -- ++(0,-0.5) -| (seq.north);
    \draw[arrow] (array.south) -- (par.north);
    \draw[arrow] (array.south) -- ++(0,-0.5) -| (chan.north);
    
    % Results
    \node[rectangle, draw, fill=white] at (-4,-1.5) {Time: 111 ms};
    \node[rectangle, draw, fill=white] at (0,-1.5) {Time: 68 ms};
    \node[rectangle, draw, fill=white] at (4,-1.5) {Time: 44 ms};
\end{tikzpicture}
\caption{Architecture Overview - Array Sum Comparison}
\end{figure}

\subsection{Kết quả thực thi}

\begin{Verbatim}[frame=single, fontsize=\small]
=== Array Sum Algorithm Comparison ===

Array size: 100,000,000
Array: [1, 2, 3, ..., 100000000]
Expected sum (formula): 5,000,000,050,000,000

============================================================

1. SEQUENTIAL SUM
----------------------------------------
   Result: 5,000,000,050,000,000
   Time: 111 ms
   Correct: True

2. PARALLEL SUM (with local state)
----------------------------------------
   Result: 5,000,000,050,000,000
   Time: 68 ms
   Correct: True

3. CHANNEL SUM (Producer-Consumer)
----------------------------------------
   Result: 5,000,000,050,000,000
   Time: 44 ms
   Correct: True

============================================================
PERFORMANCE SUMMARY
============================================================
Algorithm                 Time (ms)       Speedup
--------------------------------------------------
Sequential                111             1.00x
Parallel                  68              1.63x
Channel                   44              2.52x
============================================================
\end{Verbatim}

\subsection{Phân tích Speedup}

\begin{table}[h]
\centering
\begin{tabular}{|l|c|c|}
\hline
\textbf{Algorithm} & \textbf{Time (ms)} & \textbf{Speedup} \\
\hline
Sequential & 111 & 1.00x (baseline) \\
\hline
Parallel & 68 & 1.63x \\
\hline
Channel & 44 & \textbf{2.52x} \\
\hline
\end{tabular}
\caption{Performance Summary - Array Sum}
\end{table}

%==============================================================================
\section{Đánh giá và So sánh}
%==============================================================================

\subsection{Bảng tổng hợp các phương pháp}

\begin{table}[h]
\centering
\footnotesize
\begin{tabular}{|l|p{2.5cm}|p{2.5cm}|p{3cm}|}
\hline
\textbf{Phương pháp} & \textbf{Ưu điểm} & \textbf{Nhược điểm} & \textbf{Khi nào sử dụng} \\
\hline
Sequential & Đơn giản, dễ debug, cache-friendly & Không tận dụng đa nhân & Dữ liệu nhỏ, logic phức tạp \\
\hline
Parallel.For & Tự động phân chia công việc & Overhead với dữ liệu nhỏ & Tính toán độc lập \\
\hline
Interlocked & Thread-safe, atomic & Contention cao & Cập nhật ít, đơn giản \\
\hline
Local State & Giảm contention, hiệu quả & Code phức tạp hơn & \textbf{Reduction pattern} \\
\hline
Partitioning & Kiểm soát hoàn toàn & Phải tự quản lý & Cần tối ưu cao \\
\hline
PLINQ & Cú pháp đẹp, declarative & Overhead & Prototype, đọc dễ \\
\hline
Channel & Pipeline, decoupling & Phức tạp, overhead & Producer-Consumer, streaming \\
\hline
\end{tabular}
\caption{Tổng hợp các phương pháp lập trình song song}
\end{table}

\subsection{Khi nào nên song song hóa?}

\begin{figure}[h]
\centering
\begin{tikzpicture}[
    decision/.style={diamond, draw, minimum width=2.5cm, minimum height=1cm, align=center, inner sep=1pt},
    block/.style={rectangle, draw, minimum width=2cm, minimum height=0.8cm, align=center},
    arrow/.style={->, >=stealth},
    scale=0.85, transform shape
]
    % Start
    \node[decision, fill=yellow!20] (d1) at (0,0) {Data size\\> 10K?};
    
    % Yes branch
    \node[decision, fill=yellow!20] (d2) at (-3,-2) {Task\\độc lập?};
    \node[block, fill=green!20] (seq1) at (3,-2) {Sequential};
    
    % From d2
    \node[block, fill=blue!20] (par) at (-5,-4) {PARALLEL\\Parallel.For\\PLINQ};
    \node[decision, fill=yellow!20] (d3) at (-1,-4) {Có thể chia\\chunks?};
    
    % From d3
    \node[block, fill=orange!20] (chan) at (-3,-6) {CHANNEL\\Pipeline};
    \node[block, fill=green!20] (seq2) at (1,-6) {Sequential\\or Rethink};
    
    % Arrows
    \draw[arrow] (d1) -- node[left] {Yes} (d2);
    \draw[arrow] (d1) -- node[above] {No} (seq1);
    \draw[arrow] (d2) -- node[left] {Yes} (par);
    \draw[arrow] (d2) -- node[above] {No} (d3);
    \draw[arrow] (d3) -- node[left] {Yes} (chan);
    \draw[arrow] (d3) -- node[above] {No} (seq2);
\end{tikzpicture}
\caption{Decision Tree: Parallelize or Not?}
\end{figure}

\subsection{Best Practices}

\begin{enumerate}
    \item \textbf{Measure First}: Luôn đo đạc trước khi tối ưu
    \item \textbf{Start Simple}: Bắt đầu với Parallel.For hoặc PLINQ
    \item \textbf{Reduce Contention}: Sử dụng Local State thay vì Interlocked
    \item \textbf{Mind the Overhead}: Parallel không phải lúc nào cũng nhanh hơn
    \item \textbf{Consider Data Size}: Parallel hiệu quả với dữ liệu lớn (>10K phần tử)
\end{enumerate}

%==============================================================================
\section{Kết luận}
%==============================================================================

\subsection{Những điểm chính đã học}

\begin{enumerate}
    \item \textbf{Reduction Pattern} là pattern quan trọng trong parallel programming để tổng hợp kết quả từ nhiều thread.
    
    \item \textbf{Parallel.For với Local State} là cách tiếp cận hiệu quả nhất cho bài toán reduction, vì:
    \begin{itemize}
        \item Giảm thiểu contention giữa các thread
        \item Mỗi thread làm việc với bản sao local
        \item Chỉ đồng bộ một lần ở cuối
    \end{itemize}
    
    \item \textbf{Channel} phù hợp cho:
    \begin{itemize}
        \item Producer-Consumer pattern
        \item Pipeline processing
        \item Khi cần decouple các stage xử lý
    \end{itemize}
    
    \item \textbf{Không phải lúc nào parallel cũng nhanh hơn}:
    \begin{itemize}
        \item Overhead của thread management
        \item Cache effects
        \item Contention
    \end{itemize}
\end{enumerate}

\subsection{Recommendations}

\begin{table}[h]
\centering
\begin{tabular}{|l|l|}
\hline
\textbf{Scenario} & \textbf{Recommended Approach} \\
\hline
Simple reduction & \texttt{Parallel.For} with Local State \\
\hline
Stream processing & \texttt{Channel<T>} \\
\hline
Query-like operations & \texttt{PLINQ} \\
\hline
Fine-grained control & Manual Partitioning \\
\hline
Small data & Sequential \\
\hline
\end{tabular}
\caption{Recommendations theo từng scenario}
\end{table}

%==============================================================================
\section{Tài liệu tham khảo}
%==============================================================================

\begin{enumerate}
    \item Microsoft Docs - Parallel Programming in .NET\\
    \url{https://docs.microsoft.com/en-us/dotnet/standard/parallel-programming/}
    
    \item Microsoft Docs - System.Threading.Channels\\
    \url{https://docs.microsoft.com/en-us/dotnet/core/extensions/channels}
    
    \item Microsoft Docs - Parallel.For Method\\
    \url{https://docs.microsoft.com/en-us/dotnet/api/system.threading.tasks.parallel.for}
    
    \item Microsoft Docs - PLINQ (Parallel LINQ)\\
    \url{https://docs.microsoft.com/en-us/dotnet/standard/parallel-programming/introduction-to-plinq}
    
    \item Microsoft Docs - Thread-Safe Collections\\
    \url{https://docs.microsoft.com/en-us/dotnet/standard/collections/thread-safe/}
    
    \item Stephen Toub - Patterns of Parallel Programming\\
    \url{https://www.microsoft.com/en-us/download/details.aspx?id=19222}
    
    \item .NET Blog - An Introduction to System.Threading.Channels\\
    \url{https://devblogs.microsoft.com/dotnet/an-introduction-to-system-threading-channels/}
\end{enumerate}

\vspace{1cm}
\hrule
\vspace{0.5cm}
\begin{center}
\textit{Báo cáo được tạo ngày: 03/02/2026}\\
\textit{Môi trường thực thi: .NET 8.0, Windows, 12 CPU cores}
\end{center}

\end{document}
